\odiseachapter{marcharse}{Marcharse de Ítaca}\label{cap:marcharse}

El verano siguiente al de mi lectura de la \emph{Odisea}, lo pasé en Blanes, en la Costa Brava, con mi amigo Quelo.

Fue un verano extraño. Acabábamos de cumplir los dieciséis años. Acababa de marcharme de Alba, donde había llegado a los diez, y me había mudado con mi familia a Puerto de Sagunto, que se convertiría en mi verdadera patria chica. Pero el instituto no empezaba hasta septiembre y mi padre me dio permiso para pasar julio y agosto trabajando en el negocio que empleaba a mi amigo.

Era un invento de la España de mediados de los setenta, a la que llegaban turistas de media Europa a disfrutar de las playas y el bajo coste de la vida. Lo llevaba la Siseta, una hippy madurita y briosa, que tenía más de Circe que de Calipso. Funcionaba así: los turistas embarcaban en las golondrinas que cubrían el trayecto entre Blanes y Sant Feliu de Guíxols, deteniéndose un par de horas en Tossa de Mar. Cuando el rebaño de guiris bajaba por la pasarela desde el barco hasta el puerto, la Siseta, cámara en ristre, inmortalizaba el momento. Después salíamos disparados a revelar las fotos mientras el ganado pacía, cual vacas del sol, entre bares y tiendas de chucherías. Con los paquetes de fotos en la mano, Quelo y yo abordábamos la golondrina ---el patrón era el padre de Quelo, el tío Miquel---, localizábamos nuestras víctimas y les intentábamos colocar la foto \emph{only 25 pesetas for you, Mister}.

Nos pasábamos el día en las golondrinas. Éramos marinos, navegantes, comerciantes fenicios, generales de Cartago. Pero sobre todo, éramos aqueos camino de Ítaca, escuchando una y otra vez los versos de Kavafis en la grandiosa versión de Lluís Llach.

\begin{versocita}
Quan surts per fer el viatge cap a Ítaca,\\
has de pregar que el camí sigui llarg,\\
ple d'aventures, ple de coneixences.
\end{versocita}

No tardamos en hacernos con la traducción del griego al castellano, debida a la inspirada pluma de José María Álvarez, editada por Hiperión, la misma editorial que ha publicado la traducción de la \emph{Odisea} de Juanma Tobal.

\begin{versocita}
Si vas a emprender el viaje hacia Ítaca,\\
pide que tu camino sea largo,\\
rico en aventuras, en conocimiento.
\end{versocita}

Muchas veces he pensado que mi nave zarpó de Ítaca ese verano, dejando atrás mi niñez, llevándome por océanos que ya he recorrido por más de medio siglo. Como para Odiseo, Ítaca representa muchas cosas, algunas contradictorias. Hay una Ítaca a la que nunca he querido volver, pero aun así conserva un espacio de luz y dolor en mi corazón. Se trata de Alba, el pueblo donde campeaban cíclopes y lestrigones, pero donde encontré a Quelo, mi amigo del alma, cuyas cenizas reposan en su cementerio desde hace cinco años. Hay otra Ítaca imaginada, que nunca fue mía, la que da al Mar Menor, donde leí por primera vez a Homero, el Mar Menor que mi primo y yo surcábamos en nuestro frágil bote de madera, el Mar Menor frente al que se sentaban a jugar al dominó, las tardes de verano, mi padre y mis tíos. Y hay también la Ítaca que tenía un alto horno en lugar de castillo, la Ítaca que me recibió con una mano extendida cuando, a los dieciséis años, empezaba sexto de bachillerato. ¿Fue Atenea, siguiendo su travesti costumbre de materializarse en hombres, quien se disfrazó de muchacho, demasiado hermoso para ser humano? Recuerdo verlo venir hacia mí, el primer día de clase, recuerdo cómo buscaba yo vías de escape del aula, Odiseo experimentado en huir del bullying, recuerdo que me dijo su nombre, Juan Carlos, y yo quizás le contesté que el mío era Nadie.

Hay un momento en la vida en que la amistad es fácil, generosa y dura para siempre. Mis padres, recién llegados al Puerto, conocían de refilón a un matrimonio que tenía un hijo de mi edad. Su nombre era Toni y, como Juan Carlos, como Quelo, pasó a sumarse a la lista de mis hermanos. Hay un momento en la vida en que la patria chica lo es todo, aunque quizás aquella Ítaca industrial de mi adolescencia no era sino el país de los feacios, donde hallé la hospitalidad que no conocía aún, donde conocí no a una, sino a muchas Nausícaas, donde, más adelante, encontré y luego abandoné a Penélope, rompiendo su corazón y el mío.

En algún momento de ese largo viaje, hace ya más de treinta y cinco años, arribé a la isla de la hechicera de las largas trenzas. Y ya no pude y no quise marcharme.

Quizás por eso, en el ajuste de cuentas que tenía pendiente con el anciano Ulises, le doy, en último término, una forma de redimirse. Homero, como hemos visto, quiere mandarlo a expiar sus pecados con rituales y oraciones. Yo le propongo una alternativa mejor.

\begin{versopropio}
Hoy has decidido\\
volver a navegar.
\stanzabreak
¡Sí! ¡Al mar! Dejando atrás, esta vez para siempre,\\
tu amada isla.
\stanzabreak
¡Oh, Ítaca! Tardaste tanto tiempo en volver,\\
demasiado quizás. No encontraste lo que habías imaginado.\\
¿Y cómo podrías? Nada sobrevive sin cambiar\\
durante veinte años.
\stanzabreak
¿Qué esperabas? Habías dejado a una joven esposa\\
y a un bebé, sin edad todavía para caminar,\\
habías dejado a tus padres sanos, tu perro era solo un cachorro,\\
tu palacio estaba reluciente y bien cuidado.
\stanzabreak
Es cierto, podías culpar a los pretendientes\\
por arrasar tu hacienda. ¿Pero a quién le importa?\\
Les hiciste pagar caro sus excesos,\\
aunque tu venganza no te dejó más que manos ensangrentadas\\
y un corazón vacío. Y además, Penélope,\\
aún hermosa y tan agradable como siempre,\\
era una perfecta desconocida. Tu madre se había ido,\\
tu padre y tu perro pronto la seguirían, y tu hijo,\\
tu querido hijo, no había de quedarse.
\stanzabreak
Han pasado veinte años más desde tu regreso.\\
Según todas las normas deberías estar muerto,\\
como todos tus amigos\\
y tus enemigos\\
---en tus recuerdos\\
ya no distingues entre ellos---.\\
Y, de hecho, todos cuantos te rodean\\
están pacientemente esperando\\
que mueras. Tal vez solo tu hijo se preocupa por ti,\\
o al menos, desesperadamente,\\
tratas de creerlo.
\stanzabreak
Un anciano, que ha visto más de setenta inviernos\\
y aún, irracionalmente, se aferra a vivir.\\
Sorprendentemente fuerte, a pesar de tus dolores y enfermedades,\\
sorprendentemente lúcido, a pesar del vino oscuro\\
que te ayuda a encontrar el sueño cada noche.
\end{versopropio}

El anciano Odiseo no teme a la muerte, quizás porque, como Epicuro, ha entendido que los dioses no se ocupan de los hombres y el Hades, ese sitio horrible, está vacío de almas.

\begin{versopropio}
Es cierto que podrían elegirse otras salidas,\\
un pequeño corte en la muñeca sería suficiente\\
para enviarte, sin dolor, al inframundo,\\
o podrías mezclar en tu vino algunos de los polvos\\
que los sacerdotes utilizan\\
para facilitar el paso de los moribundos.
\stanzabreak
Aunque, pensándolo bien, si tuvieses que elegir,\\
probablemente sería el ir nadando\\
hacia la línea del horizonte,\\
hacia la isla de Circe.
\end{versopropio}

Pero, sobre todo, este anciano entiende, por fin, lo que significa amar.

\begin{versopropio}
Circe,\\
qué extraño, finalmente, admitir\\
que la amabas.\\
La amas todavía, con toda tu alma,\\
escapaste de su isla por orgullo\\
y por miedo.
\stanzabreak
Tú, que sabías que nunca\\
se inclinaría ante tus deseos,\\
que no se sometería a tus caprichos,\\
que no cedería por tu bien,\\
ni la adulación ni las amenazas\\
la ablandarían.\\
Ella te conocía por lo que eras,\\
un mercenario, un villano y un matón,\\
podía verte por dentro, sabía\\
lo oscura que era tu alma.
\end{versopropio}

Extraño que la hechicera, a quien no podía engañar ni manipular, lo amara de todas formas.

\begin{versopropio}
Irónicamente,\\
sin embargo, ella te amaba,\\
pero sus ojos eran un espejo\\
donde veías, cada día, tu verdadero rostro.\\
¿Y qué veías detrás de la fuerte mandíbula,\\
la nariz viril, la barbilla firme?\\
Veías un monstruo\\
que no merecía ser amado.
\stanzabreak
En consecuencia, huiste,\\
diciendo a todo el mundo lo feliz que eras,\\
lo que te costaba escapar.\\
Una mentira más, por supuesto.\\
No intentó retenerte, no se quejó\\
ni pidió explicaciones. Era Circe, orgullosa,\\
hermosa e indomable.
\end{versopropio}

Extraño darse cuenta, al final de la vida, de que el final todavía no ha llegado.

\begin{versopropio}
Ahora ha llegado el momento\\
para que intentes encontrar tu camino\\
de regreso a su isla.\\
¿Seguirá viva?\\
---Diosa o no, sabes muy bien\\
que nadie vive para siempre---.\\
¿Se acordará aún de ti?\\
¿Seguirás importándole?\\
No importa. Volverás a sus costas,\\
desnudo. Como cuando llegaste a la isla de Nausícaa.\\
Pero esta vez no llevarás ni ropa\\
ni máscaras. No ocultarás tu culpa\\
ni tus pecados. No mentirás,\\
no pedirás que te acepte\\
o te perdone. Pero si ella te permite\\
habitar en su casa,\\
dormir en una alfombra cerca de su cama,\\
caminar a su lado,\\
tú te darás por satisfecho.
\end{versopropio}

Extraño comprender que le queda una misión por delante. Encontrar a Circe y liberar a todos los fantasmas que le atormentan.

\begin{versopropio}
Tal vez todavía ella tenga la droga\\
que una vez te dejó caminar por el Hades\\
---quizás, has llegado a pensar,\\
el Hades no es más que las profundidades\\
de tu mente atormentada---.\\
La tomarás de nuevo, así soltando\\
a todos los fantasmas allí encadenados,\\
liberando a todos esos inocentes,\\
dejando el lugar, para siempre, vacío.
\end{versopropio}

Y así, finalmente, Odiseo vuelve a embarcarse:

\begin{versopropio}
Hoy navegas de nuevo.\\
Tu hijo ha ofrecido uno de sus barcos\\
y una tripulación experta. Aún mejor,\\
se ofreció a navegar contigo.\\
No es que todo el oro de la Tierra\\
ni todos los diamantes y zafiros\\
ni todos los collares de perlas y rubíes ni las ricas sedas\\
que arrancaste de Troya,\\
de la ciudad de los cícones\\
o de las cuevas de los cíclopes\\
fueran suficientes para pagar el pasaje de Telémaco.\\
Te lanzarías a las olas\\
antes de dejarle navegar contigo.\\
Pertenece a su familia, se quedará\\
y será el padre que tú nunca fuiste.\\
Pero ¡qué dulce su oferta en tu corazón,\\
como la miel en los besos de Penélope,\\
hace ya tantas primaveras!
\end{versopropio}

Marchándose de Ítaca, porque el destino de los héroes ---y todos los seres humanos lo somos--- es abandonar lo que más se ama, para poder seguir amando.

\begin{versopropio}
¡Sí! ¡Al mar!\\
dejando atrás el miedo,\\
dejando los remordimientos,\\
dejando el hogar.\\
Para encontrar el verdadero amor,\\
para encontrar a Circe.
\end{versopropio}
